\chapter[ALGORITMOS DE OTIMIZAÇÃO PARA POLÍTICAS SUPERVISÓRIAS]{ALGORITMOS DE OTIMIZAÇÃO PARA POLÍTICAS SUPERVISÓRIAS}
\label{apendice:otimizacao}

% TODO: desenvolver este apêndice. Esboço da ideia central, a ser detalhado:
%
% A camada de RTO (Seção~\ref{sec:plantwide}) requer um procedimento explícito para
% transformar a função de custo J em política supervisória, isto é, em setpoints
% concretos para as variáveis primárias c. Este apêndice explora algoritmos de
% otimização candidatos a desempenhar esse papel -- o "cérebro" das políticas do
% supervisor Kubernetes -- partindo de funções objetivo definidas a priori,
% executando a otimização sobre cada uma e registrando os resultados como
% políticas supervisórias disponíveis.
%
% A função de custo do TEP (Eq.~\ref{eq:tep_cost}) pode ser estendida para
% incorporar outras dimensões de desempenho operacional além das parcelas
% originais de Downs e Vogel, compondo um objetivo composto:

\begin{equation}
	J =
	J_{\text{operação}}
	+ J_{\text{qualidade}}
	+ J_{\text{energia}}
	+ J_{\text{risco}}
	+ J_{\text{manutenção}}
	+ J_{\text{produção perdida}}
	\label{eq:tep_cost_extended}
\end{equation}

% TODO: detalhar cada parcela, os algoritmos de otimização avaliados e como o
% resultado de cada otimização se traduz em política supervisória.

% -------------------------------------------------------------------------
\section{Algoritmos por tipo de modelo}
\label{sec:apendice_algoritmos}

% TODO (rascunho a ser detalhado): esta seção apenas situa, de forma tosca, qual
% família de algoritmo resolve qual tipo de problema. Os exemplos do próprio
% \citeonline{Williams2013} já são resolvidos por programação matemática: o livro
% aponta três famílias principais de algoritmos usadas por pacotes comerciais --
% simplex revisado para programação linear, extensão separável do simplex para
% programação separável, e branch-and-bound para programação inteira.

Para os casos relevantes a este trabalho, o tipo de modelo resultante da função de
custo $J$ (Equação~\ref{eq:tep_cost_extended}) determina a família de algoritmo
candidata, conforme a Tabela~\ref{tab:algoritmos_otimizacao}.

\begin{table}[ht!]
\centering
\caption{Famílias de algoritmos por tipo de modelo (rascunho)}
\label{tab:algoritmos_otimizacao}
\footnotesize
\begin{tabular}{|p{3.0cm}|p{3.3cm}|p{6.0cm}|}
\hline
\textbf{Caso} & \textbf{Tipo de modelo} & \textbf{Algoritmo típico} \\
\hline
Blending & LP & Simplex / interior-point \\
\hline
Refinery optimization & LP, ou NLP se houver não linearidade & Simplex / interior-point; eventualmente NLP \\
\hline
Factory Planning 1 & LP multi-período & Simplex / interior-point \\
\hline
Factory Planning 2, com manutenção & MILP & Branch-and-bound / branch-and-cut \\
\hline
Scheduling, lógica, escolhas discretas & MILP / CP-SAT & Branch-and-bound, branch-and-cut, constraint programming \\
\hline
MPC linear & QP & Active-set, interior-point, ADMM \\
\hline
MPC não linear / TEP completo & NLP / NMPC & SQP, interior-point \\
\hline
Planejamento com incerteza & Stochastic programming & LP/MILP com cenários, decomposição \\
\hline
\end{tabular}
\fonte{Elaborado pelo autor, com base em \citeonline{Williams2013}.}
\end{table}

% TODO: para cada linha relevante ao TEP (em especial MPC não linear / NMPC, que é
% o caso do problema completo de plant-wide control), detalhar o algoritmo
% escolhido, sua formulação e como o resultado se traduz em política
% supervisória concreta para o Kubernetes.
